% SPDX-License-Identifier: Apache-2.0
%
% TxHDL from zero: a blinky in Bazel, from an empty directory to a
% netlist, for somebody who has neither the tree nor Bazel. Built by
% //docs:zero. Every file shown is included from tutorial/blinky/,
% which the tree builds as a check, and every printout is what the
% build produced by running it.
\documentclass[journal,9pt]{IEEEtran}

\newcommand{\listingsize}{\fontsize{6}{7}\selectfont}
% SPDX-License-Identifier: Apache-2.0
%
% The house preamble, shared by every document in this directory. Loaded
% after \documentclass, before \title. Keeping it in one file is what
% stops two articles drifting apart in font, listing style or figure
% style.
% Computer Modern. IEEEtran selects Times (ptm) by default, so the face has
% to be chosen explicitly.
%
% Latin Modern is the Computer Modern variety used here, and the choice is
% forced rather than aesthetic. Under T1 encoding, plain `cmr` has no Type 1
% outlines unless cm-super is installed, and it is not in the pinned
% distribution, so pdflatex falls back to EC bitmaps and every glyph in the
% PDF becomes a Type 3 font. That was measured: the first build produced a
% document with 10 Type 3 fonts and no Type 1. Latin Modern is Computer
% Modern redrawn with a full T1 repertoire and real .pfb outlines, so the
% same design comes out scalable.
\usepackage[T1]{fontenc}
\usepackage[utf8]{inputenc}
\usepackage{lmodern}
% microtype is not in the pinned TeX distribution. emergencystretch alone
% keeps the two-column measure from producing overfull lines.
\setlength{\emergencystretch}{3em}

\usepackage{amsmath}
\usepackage{amssymb}
\usepackage{amsthm}
\usepackage{booktabs}
\usepackage{array}
% enumitem is not in the pinned TeX distribution; the list spacing below
% does what \setlist would have done.
\usepackage{float}
\usepackage{xcolor}
\usepackage{listings}
\usepackage{tikz}
\usetikzlibrary{arrows.meta,positioning,fit,backgrounds,calc}
\usepackage[hidelinks,bookmarks=true]{hyperref}

% Numbered, definition-style box for a rule the reader is meant to apply.
\theoremstyle{definition}
\newtheorem{rulebox}{Rule}
\newtheorem{example}{Example}

% Unnumbered asides.
\theoremstyle{remark}
\newtheorem*{tip}{Tip}
\newtheorem*{pitfall}{Pitfall}

% Tighter lists, in place of enumitem's \setlist.
\setlength{\topsep}{2pt}
\setlength{\itemsep}{1pt}
\setlength{\parsep}{0pt}

\newcommand{\code}[1]{\texttt{#1}}
\newcommand{\term}[1]{\emph{#1}}

% Syntax highlighting. Four roles, distinguishable in colour and still
% distinguishable in greyscale, because an IEEE article gets printed: the
% keyword is the only bold role, the comment is the only italic one, and the
% two remaining roles differ in lightness as well as in hue.
\definecolor{kwblue}{RGB}{20,60,150}
\definecolor{cmtgreen}{RGB}{90,120,90}
\definecolor{strred}{RGB}{150,50,40}
\definecolor{numgray}{gray}{0.45}
\definecolor{framegray}{gray}{0.70}
\definecolor{bgshade}{gray}{0.97}
% A darker fill for figure cells. The listing background has to stay very
% light to keep code readable; a figure cell has to be clearly shaded or the
% distinction it draws is invisible in print.
\definecolor{cellshade}{gray}{0.90}

% The listing size is the document's choice, because it depends on the
% column: a two-column article sets `\listingsize` to 6pt and keeps its
% listings within 70 columns, a one-column document keeps the body size
% and includes source files at 80. Lines are never broken; a line that does not fit
% overflows the frame, which is visible, rather than wrapping, which is
% not.
\providecommand{\listingsize}{\scriptsize}

% `'` is deliberately not a string delimiter: the language uses it for loop
% labels, as Rust does, so treating it as a quote would swallow the rest of
% the line.
\lstdefinestyle{txhdl}{
  basicstyle=\ttfamily\listingsize,
  keywordstyle=\bfseries\color{kwblue},
  commentstyle=\itshape\color{cmtgreen},
  stringstyle=\color{strred},
  numberstyle=\tiny\color{numgray},
  morekeywords={mod,use,pub,const,type,struct,enum,trait,impl,fn,async,let,mut,
    match,if,else,loop,while,for,in,break,continue,return,as,await,
    module,bus,channel,transaction,protocol,pipeline,flow,seq,config,
    port,stage,pipe,instance,connect,bind,tag,domain,blackbox,foreign,
    property,role,signal,handler,true,false,
    sequence,interface,modport,implements,package,import,var,wire,func,proc,
    case,default,next,fork,branch,join,state,fsm},
  morecomment=[l]{//},
  morestring=[b]",
  showstringspaces=false,
  columns=fullflexible,
  keepspaces=true,
  breaklines=false,
  % Framed, so a listing is bounded rather than floating in the column.
  frame=single,
  rulecolor=\color{framegray},
  framesep=3pt,
  backgroundcolor=\color{bgshade},
  xleftmargin=3pt,
  xrightmargin=3pt,
  aboveskip=6pt,
  belowskip=6pt,
  % Every listing is numbered and captioned, so it can be referred to by
  % number. `caption` without `float` numbers the listing and leaves it
  % where it was written, which is what a listing in running prose needs.
  captionpos=b,
  abovecaptionskip=2pt,
  belowcaptionskip=6pt,
}

% Figures drawn as text. Framed the same way, and with the highlighting
% switched off, because a box-and-arrow diagram is not code and half its
% words turning blue reads as an accident. Setting a style adds keys rather
% than clearing the ones already set, so the three roles are neutralised by
% hand rather than by leaving them out.
\lstdefinestyle{ascii}{
  basicstyle=\ttfamily\listingsize,
  keywordstyle=\ttfamily,
  commentstyle=\ttfamily,
  stringstyle=\ttfamily,
  columns=fullflexible,
  keepspaces=true,
  frame=single,
  rulecolor=\color{framegray},
  framesep=3pt,
  backgroundcolor=\color{bgshade},
  xleftmargin=3pt,
  xrightmargin=3pt,
  aboveskip=6pt,
  belowskip=6pt,
}
\lstset{style=txhdl}

% House TikZ styles. `shorten` lives on the arrow styles rather than on each
% arrow, so every arrow in the document gets the gap, including ones added
% later. TikZ clips a path at a node's border, so without it an arrowhead
% butts against the box with no gap at all. Keep `node distance` well above
% twice the shorten, or the arrow shrinks to nothing.
\tikzset{
  box/.style={draw,rounded corners=2pt,align=center,inner sep=4pt,
              font=\footnotesize},
  boxf/.style={box,fill=bgshade},
  lbl/.style={font=\scriptsize\itshape,align=center},
  fl/.style={-{Stealth[length=4pt]},shorten >=2.5pt,shorten <=2.5pt},
  fld/.style={fl,dashed},
}



\InputIfFileExists{buildstamp.tex}{}{}
\providecommand{\buildstamp}{Draft build (outside the Bazel flow).}

\title{TxHDL from Zero: A Blinky in Bazel}

\author{Filip Filmar and Dragiša Janković%
\thanks{Both authors are independent. Filip Filmar is the
corresponding author, at f@hdlfactory.com; Dragiša Janković is
reached at linkedin.com/in/dragisa-jankovic.}%
\thanks{This is the tutorial for an empty directory: nothing
installed, no copy of the TxHDL tree, and a blinky built, run and
lowered to Verilog with Bazel at the end. Every file it shows is
included from the project repository \code{HDL/txhdl}, under
\code{tutorial/blinky/}, which the project's own build builds as a
check, and every printout is what the build produced by running it.
The language is learnt from the step-by-step tutorial~\cite{tutorial}
and the cheat sheet~\cite{cheatsheet}; the cover~\cite{cover} lists
every document.}%
\thanks{The authors used a large language model, Claude, as an
assistant in exploring the concepts and in writing and constructing
the documents and the programs; every commit of the repository says
so and carries the prompts that led to it, verbatim.}%
\thanks{\buildstamp}}

\markboth{TxHDL from Zero}{Filmar and Janković: TxHDL from Zero}

\begin{document}
\maketitle

\begin{abstract}
TxHDL is a hardware description language that is a Rust library, and
a design in it is built by Bazel. This document starts from an empty
directory and ends with a blinky that runs, prints its LED as a line
of characters, and writes its Verilog as a file of the build. It
takes four files and one binary on the machine.
\end{abstract}

\section{What You Need}
\label{sec:z-need}

One program: Bazelisk, a single binary from its GitHub releases page
or a package manager, put on the path as \code{bazel}. It reads the
version a workspace names and fetches that Bazel. Everything else
comes with the build: the Rust compiler, fetched by the Rust rules at
the version TxHDL pins; TxHDL itself, fetched from its GitHub mirror
at a named commit; and TxHDL's own dependencies, resolved through two
registries. Nothing is installed on the machine, and no Verilog tool
is needed to get a netlist, since the netlist is what the Rust
program prints. Git is needed once, for the fetch of the mirror.

\section{Four Files}
\label{sec:z-files}

A workspace is a directory with these files in it. The first names
the Bazel version, so that Bazelisk fetches the one this document was
built with:

\lstinputlisting[style=ascii,caption={\code{.bazelversion}: one line.},label={lst:z-version}]{tutorial/blinky/.bazelversion}

The second names the registries Bazel looks modules up in. TxHDL's
dependencies, the Rust rules among them, come from the Bazel Central
Registry; a few of its tools come from the project's own registry,
named first:

\lstinputlisting[style=ascii,caption={\code{.bazelrc}: where modules come from.},label={lst:z-rc}]{tutorial/blinky/.bazelrc}

The third is the module: a name for the workspace, the Rust rules,
and TxHDL, which no registry holds yet, so it is named by its
repository and a commit. The commit here is the one this document
was built from; a newer one is any commit of the mirror's
\code{main}, and the version, \code{0.0.0}, is what TxHDL calls
itself.

\lstinputlisting[style=ascii,caption={\code{MODULE.bazel}: two dependencies, one of them by commit.},label={lst:z-module}]{tutorial/blinky/MODULE.bazel}

The fourth is the build file. The blinky is a Rust binary that
depends on TxHDL's library and its macros; the second rule runs it
with \code{--verilog} and keeps what it prints as a file of the
build, \code{blinky.v}. The visibility line and the last rule are
for the document you are reading, which runs the blinky from the
project's tree, and are not needed in a workspace of your own.

\lstinputlisting[style=ascii,caption={\code{BUILD.bazel}: the binary, its netlist, and the sources.},label={lst:z-build}]{tutorial/blinky/BUILD.bazel}

\section{The Blinky}
\label{sec:z-blinky}

The design is one register and one output. A unit is a struct of its
state, here the count, and its behaviour is the \code{run} of the
\code{Unit} trait: one loop, whose first statement waits for the
rising edge of the clock, so that each iteration is one cycle. The
count wraps at the period and counts otherwise, which \code{when!}
says with the condition first; the LED is a compare on the count, a
wire, driven with \code{set}. The attribute \code{\#[lower]} reads
the loop and writes \code{Blinky::verilog} beside it, so the same
struct that simulates also prints its netlist. The \code{main} runs
the design for two periods and prints the LED, or prints the Verilog
when asked.

\lstinputlisting[caption={\code{blinky.rs}: the design, and a \code{main} that runs it or prints it.},label={lst:z-rs}]{tutorial/blinky/blinky.rs}

\section{Build and Run}
\label{sec:z-run}

In the directory:

\begin{lstlisting}[style=ascii]
bazel run //:blinky
\end{lstlisting}

The first run fetches Bazel, the Rust toolchain and TxHDL with its
dependencies, and compiles them: two minutes on the machine this
document was built on, and some three gigabytes on disk, most of it
the Rust toolchain and the compiler's outputs; every run after it is
seconds, since Bazel keeps what it fetched and built. The run prints one character per cycle,
the LED high for the first eight cycles of each period of sixteen:

\lstinputlisting[style=ascii,caption={What \code{bazel run //:blinky} printed when this document was built.},label={lst:z-out}]{out_zero.txt}

\section{The Netlist as a File}
\label{sec:z-netlist}

The second rule of the build file keeps the Verilog:

\begin{lstlisting}[style=ascii]
bazel build //:blinky.v
cat bazel-bin/blinky.v
\end{lstlisting}

It is the module below: the register, the wrap as an \code{if} in
the clocked block, the LED as a continuous assignment. The constants
of the source, \code{PERIOD - 1} and \code{PERIOD / 2}, are the
numbers they are, since the lowering folds them. The file is what any
Verilog tool takes; the Vreteno document~\cite{vreteno} shows the
whole flow to a board, with Vivado installed by the build.

\lstinputlisting[style=ascii,caption={\code{blinky.v}, as the build wrote it.},label={lst:z-v}]{tutorial/blinky/blinky.v}

\section{From Here}
\label{sec:z-next}

Change \code{PERIOD} and run again; the line and the netlist follow.
Add a second register and the trace derive names it in a waveform,
which the step-by-step tutorial~\cite{tutorial} shows how to write and
draw, along with channels between units and the checks the project
runs on every netlist against its own run. The cheat
sheet~\cite{cheatsheet} is the language on one page, and the examples
document~\cite{examples} has one example per concept, each with its
printout and its waveform. To move the pin to a commit of your own
choosing, change the commit in \code{MODULE.bazel} and run again;
Bazel fetches the difference.

\begin{thebibliography}{9}

\bibitem{tutorial}
F.~Filmar and D.~Janković, ``TxHDL, step by step,'' project
repository \code{HDL/txhdl}, \code{//docs:tutorial}, 2026.

\bibitem{cheatsheet}
F.~Filmar and D.~Janković, ``TxHDL cheat sheet,'' project repository
\code{HDL/txhdl}, \code{//docs:cheatsheet}, 2026.

\bibitem{examples}
F.~Filmar and D.~Janković, ``TxHDL by example,'' project repository
\code{HDL/txhdl}, \code{//docs:examples}, 2026.

\bibitem{vreteno}
F.~Filmar and D.~Janković, ``Vreteno: an RV32IM core in TxHDL,''
project repository \code{HDL/txhdl}, \code{//docs:vreteno}, 2026.

\bibitem{cover}
F.~Filmar and D.~Janković, ``TxHDL documents,'' project repository
\code{HDL/txhdl}, \code{//docs:cover}, 2026.

\end{thebibliography}

\end{document}
